\documentclass[aspectratio=169,12pt]{beamer}
\usepackage{amsmath}
\usepackage{amssymb}
\usepackage{array}
\usepackage[most]{tcolorbox}
\usepackage{graphicx}
\usepackage{tikz}
\usetikzlibrary{shadings}
\usepackage[sfdefault,lf]{FiraSans}
\renewcommand{\familydefault}{\sfdefault}
\setlength{\parindent}{0pt}
\everymath{\displaystyle}

%% TikZ 3D shape macros for surface-area notes
% DrawCuboid[3]{depth}{width}{height}
\newcommand{\DrawCuboid}[3]{%
\begin{tikzpicture}[line join=round,line cap=round]
  \pgfmathsetmacro{\ddx}{0.18*#1}
  \pgfmathsetmacro{\ddy}{0.08*#1}
  \pgfmathsetmacro{\wx}{0.22*#2}
  \pgfmathsetmacro{\hy}{0.28*#3}
  \coordinate (o) at (0,0);
  \coordinate (a) at (\wx,0);
  \coordinate (b) at (\wx,\hy);
  \coordinate (c) at (0,\hy);
  \coordinate (d) at (\ddx,\ddy);
  \coordinate (e) at (\wx+\ddx,\ddy);
  \coordinate (f) at (\wx+\ddx,\hy+\ddy);
  \coordinate (g) at (\ddx,\hy+\ddy);
  \fill[mmgreenLeaf!15] (o) -- (a) -- (b) -- (c) -- cycle;
  \fill[mmgreenLeaf!25] (a) -- (e) -- (f) -- (b) -- cycle;
  \fill[mmgreenLeaf!10] (c) -- (b) -- (f) -- (g) -- cycle;
  \draw (a) -- (b) -- (c) -- (o) -- (a);
  \draw (b) -- (f) -- (g) -- (c);
  \draw (f) -- (e) -- (a);
  \draw[densely dashed] (o) -- (d) -- (e);
  \draw[densely dashed] (d) -- (g);
  \draw[<->,thin] ([xshift=-4pt]o) -- ([xshift=-4pt]c) node[midway,left,font=\tiny]{#3};
  \draw[<->,thin] ([yshift=-4pt]o) -- ([yshift=-4pt]a) node[midway,below,font=\tiny]{#2};
  \draw[<->,thin] ([xshift=3pt,yshift=-3pt]a) -- ([xshift=3pt,yshift=-3pt]e) node[midway,sloped,below,font=\tiny]{#1};
  \path (o) rectangle (f);
\end{tikzpicture}%
}
% DrawCubeNet[2]{side length}{label}
\newcommand{\DrawCubeNet}[2]{%
\begin{tikzpicture}[x=0.18cm,y=0.18cm,line cap=round,line join=round]
  \pgfmathsetmacro{\s}{#1}
  % Cross net: 6 squares arranged like a T
  % Center square (front)
  \fill[mmgreenLeaf!20] (0,0) rectangle (\s,\s);
  \draw (0,0) rectangle (\s,\s);
  \node[font=\tiny] at (\s/2,\s/2) {#2};
  % Top (back)
  \fill[mmgreenLeaf!10] (0,\s) rectangle (\s,2*\s);
  \draw (0,\s) rectangle (\s,2*\s);
  \node[font=\tiny] at (\s/2,1.5*\s) {#2};
  % Bottom (front)
  \fill[mmgreenLeaf!20] (0,-\s) rectangle (\s,0);
  \draw (0,-\s) rectangle (\s,0);
  \node[font=\tiny] at (\s/2,-\s/2) {#2};
  % Left (left side)
  \fill[mmgreenLeaf!15] (-\s,0) rectangle (0,\s);
  \draw (-\s,0) rectangle (0,\s);
  \node[font=\tiny] at (-\s/2,\s/2) {#2};
  % Right (right side)
  \fill[mmgreenLeaf!15] (\s,0) rectangle (2*\s,\s);
  \draw (\s,0) rectangle (2*\s,\s);
  \node[font=\tiny] at (1.5*\s,\s/2) {#2};
  % Back (top continuation)
  \fill[mmgreenLeaf!10] (0,2*\s) rectangle (\s,3*\s);
  \draw (0,2*\s) rectangle (\s,3*\s);
  \node[font=\tiny] at (\s/2,2.5*\s) {#2};
\end{tikzpicture}%
}\newcommand{\DrawCylinder}[2]{%
\begin{tikzpicture}[x=0.55cm,y=0.55cm,line cap=round,line join=round]
  \fill[top color=gray!50!black,bottom color=gray!10,middle color=gray,shading=axis,opacity=0.25]
    (0,-0.9) ellipse [x radius=0.78, y radius=0.22];
  \fill[left color=gray!50!black,right color=gray!50!black,middle color=gray!50,shading=axis,opacity=0.25]
    (0.78,-0.9) -- (0.78,0.9) arc[start angle=360,end angle=180,x radius=0.78,y radius=0.22] -- (-0.78,-0.9) arc[start angle=180,end angle=360,x radius=0.78,y radius=0.22];
  \fill[top color=gray!90!,bottom color=gray!2,middle color=gray!30,shading=axis,opacity=0.25]
    (0,0.9) ellipse [x radius=0.78, y radius=0.22];
  \draw (-0.78,0.9) -- (-0.78,-0.9) arc[start angle=180,end angle=360,x radius=0.78,y radius=0.22] -- (0.78,0.9);
  \draw (0,0.9) ellipse [x radius=0.78, y radius=0.22];
  \draw[densely dashed] (-0.78,-0.9) arc[start angle=180,end angle=0,x radius=0.78,y radius=0.22];
  \draw[<->,thin] (1.15,-0.9) -- (1.15,0.9) node[midway,right,font=\tiny]{#1};
  \draw[<->,thin] (-0.78,-1.25) -- (0.78,-1.25) node[midway,below,font=\tiny]{$2\times #2$};
\end{tikzpicture}%
}\setbeamertemplate{navigation symbols}{}
\setbeamertemplate{footline}{}
\setbeamertemplate{headline}{}
\setbeamertemplate{blocks}[rounded][shadow=false]

\definecolor{mmpageWhite}{HTML}{FCFBF7}
\definecolor{mmtanSoft}{HTML}{F7F5EF}
\definecolor{mmgreenDeep}{HTML}{244B2C}
\definecolor{mmgreenLeaf}{HTML}{5D8A56}
\definecolor{mmgreenSoft}{HTML}{E4EFD9}
\definecolor{mmtext}{HTML}{163221}
\definecolor{mmwarn}{HTML}{8F2F36}
\definecolor{mmwarnSoft}{HTML}{F8E8EA}

\setbeamercolor{background canvas}{bg=mmpageWhite}
\setbeamercolor{normal text}{fg=mmtext,bg=mmpageWhite}
\setbeamercolor{block title}{fg=mmgreenDeep,bg=mmgreenSoft}
\setbeamercolor{block body}{fg=mmtext,bg=mmtanSoft}
\setbeamercolor{block title example}{fg=mmgreenDeep,bg=mmgreenSoft}
\setbeamercolor{block body example}{fg=mmtext,bg=white}
\setbeamercolor{block title alerted}{fg=mmwarn,bg=mmwarnSoft}
\setbeamercolor{block body alerted}{fg=mmtext,bg=white}
\setbeamercolor{itemize item}{fg=mmgreenDeep}
\setbeamercolor{itemize subitem}{fg=mmgreenDeep}

\newtcolorbox{MMCard}[2][]{enhanced,breakable=false,colback=#2,colframe=black,boxrule=1.5pt,arc=8pt,left=5pt,right=5pt,top=4pt,bottom=4pt,before skip=0pt,after skip=0pt,#1}
\newcommand{\KoruIcon}{\raisebox{-0.2em}{\includegraphics[height=1.05em]{../../../manamaths/SITE/header-logo.png}}}
\newcommand{\NotesTitle}[1]{{\colorbox{mmtanSoft}{\parbox{0.975\linewidth}{\vspace{0.14em}\hspace{0.25em}{\LARGE\bfseries\textcolor{mmgreenDeep}{#1}\hfill\KoruIcon}\vspace{0.14em}}}\par\vspace{0.18em}{\color{mmgreenLeaf}\rule{\linewidth}{2.0pt}}\par\vspace{0.12em}}}
\newcommand{\SectionTitle}[1]{{\large\bfseries\textcolor{mmgreenDeep}{#1}}\par\vspace{0.04em}}
\newcommand{\GreenDivider}{\par\vspace{0.01em}{\color{mmgreenLeaf}\rule{\linewidth}{2.4pt}}\par\vspace{0.03em}}
\newcommand{\KeyIdea}[1]{\begin{MMCard}[title={\textbf{Key idea}},colbacktitle=mmgreenSoft,coltitle=mmgreenDeep]{mmtanSoft}\small #1\end{MMCard}}
\newcommand{\WorkedExample}[2]{\begin{MMCard}[title={\textbf{#1}},colbacktitle=mmgreenSoft,coltitle=mmgreenDeep]{white}\small #2\end{MMCard}}
\newcommand{\CommonMistake}[1]{\begin{MMCard}[title={\textbf{Common mistake}},colbacktitle=mmwarnSoft,coltitle=mmwarn]{white}\small #1\end{MMCard}}
\newcommand{\TryThese}[1]{\begin{MMCard}[title={\textbf{Try these}},colbacktitle=mmgreenSoft,coltitle=mmgreenDeep]{mmtanSoft}\small #1\end{MMCard}}
\newcommand{\StepLine}[2]{\textbf{#1.} #2\par\vspace{0.03em}}

\setbeamertemplate{background}{
 \begin{tikzpicture}[remember picture,overlay]
 \draw[black,line width=1.5pt,rounded corners=10pt]
 ([xshift=0.45cm,yshift=-0.45cm]current page.north west)
 rectangle
 ([xshift=-0.45cm,yshift=0.45cm]current page.south east);
 \end{tikzpicture}
}

\begin{document}

\begin{frame}[t,shrink=10]
\NotesTitle{Surface area}

\begin{columns}[T,onlytextwidth]
\begin{column}{0.48\textwidth}
\KeyIdea{Surface area is the total area of all the outside faces of a 3D shape. Imagine wrapping the shape in paper --- the amount of paper needed is the surface area. A good method is to find the area of each face, then add them all together.}

\SectionTitle{Step-by-step method}
\StepLine{1}{Identify every face on the outside of the shape.}
\StepLine{2}{Find the area of each face using the correct formula.}
\StepLine{3}{Multiply when faces are the same (congruent pairs).}
\StepLine{4}{Add all the areas together.}
\StepLine{5}{Write the final answer with square units (cm$^2$, m$^2$, etc.).}

\CommonMistake{Forgetting a face. A cube and a cuboid each have \textbf{6 faces}. A closed cylinder has \textbf{2 circles} plus the curved surface --- that's 3 parts, not 2!}
\end{column}
\begin{column}{0.48\textwidth}
\WorkedExample{Cube}{%
\begin{minipage}[t]{0.42\linewidth}
\centering\DrawCuboid{4}{4}{4}
\end{minipage}%
\hfill
\begin{minipage}[t]{0.52\linewidth}
\centering\DrawCubeNet{4}{$4\times4=16$}
\end{minipage}
\\[4pt]
A cube has side length 4~cm. All faces are identical squares.\\[4pt]
Area of one face: $4\times4 = 16$~cm$^2$\\
6 identical faces: $6\times16 = 96$~cm$^2$\\[4pt]
\GreenDivider
\textbf{Surface area = 96~cm$^2$}}

\WorkedExample{Cuboid}{%
\DrawCuboid{8}{3}{2}\\[2pt]
A box is 8~cm by 3~cm by 2~cm. Opposite faces are the same size.\\[4pt]
Top \& bottom: $2\times (8\times 3)=48$\\
Front \& back: $2\times (8\times 2)=32$\\
Left \& right: $2\times (3\times 2)=12$\\
Total: $48+32+12=92$\\[4pt]
\GreenDivider
\textbf{Surface area = 92~cm$^2$}}
\end{column}
\end{columns}
\GreenDivider
\begin{columns}[T,onlytextwidth]
\begin{column}{0.48\textwidth}
\WorkedExample{Cylinder}{%
\DrawCylinder{10}{3}\\[2pt]
A cylinder has radius 3~cm and height 10~cm. It has two circles and one curved surface (which unfolds to a rectangle).\\[4pt]
Two circles: $2\pi r^2 = 2\pi(3^2)=18\pi$\\
Curved surface: $2\pi rh = 2\pi(3)(10)=60\pi$\\
Total: $18\pi+60\pi = 78\pi$\\[4pt]
\GreenDivider
\textbf{Surface area = $78\pi$~cm$^2$}}

\CommonMistake{Using the diameter instead of the radius. For a cylinder with diameter 8~cm, the radius is 4~cm. Always halve the diameter before using the formula.}
\end{column}
\begin{column}{0.48\textwidth}
\TryThese{\begin{enumerate}
\item A cube has side length 6~cm. Find its surface area.
\item A cuboid is 10~cm by 4~cm by 2~cm. Find its surface area.
\item A cylinder has radius 4~cm and height 9~cm. Write the total surface area in terms of $\pi$.
\end{enumerate}
\vspace{4pt}
\SectionTitle{Answers}
\begin{enumerate}
\item $6\times(6\times6)=216$~cm$^2$
\item $2(10\times4)+2(10\times2)+2(4\times2)=136$~cm$^2$
\item $2\pi(4^2)+2\pi(4)(9)=32\pi+72\pi=104\pi$~cm$^2$
\end{enumerate}}
\end{column}
\end{columns}
\end{frame}

\end{document}
